\section{One-dimensional data}\label{sec:oneD}

\subsection{Pipeline overview}
The one-dimensional branch is organised as a \emph{load $\to$ process $\to$
plot $\to$ analyse} pipeline. Each stage is a module under
\code{pymorgan.oneD}, and the \code{Dataset1D} object ties them together:
it holds the raw arrays and exposes one set of methods per stage. The stage
functions take the dataset as their first argument, so they remain reusable on
their own and new routines integrate with the whole pipeline. Table~\ref{tab:stages}
maps each stage to its module and entry points.

\begin{table}[h]
  \centering
  \begin{tabular}{lll}
    \toprule
    Stage & Module & Entry point(s) \\
    \midrule
    load    & \code{pymorgan.oneD.load}    & \code{pm.load\_1D}, \code{Dataset1D.from\_file} \\
    process & \code{pymorgan.oneD.process} / \code{chirp} & \code{Dataset1D.background\_correct}, \code{apply\_chirp\_correction} \\
    plot    & \code{pymorgan.oneD.plot}    & \code{Dataset1D.plot\_*} \\
    analyse & \textit{(PyRATE-TA)} & \textit{separate project} \\
    \bottomrule
  \end{tabular}
  \caption{The 1-D pipeline stages and where they live.}
  \label{tab:stages}
\end{table}

\subsection{Data model}
A one-dimensional measurement is represented by a \code{Dataset1D} object,
which bundles the arrays produced by a loader. The signal is stored as a
three-dimensional array to generalise over multiple detectors,
\begin{equation}
  Z \in \mathbb{R}^{N_\tau \times N_\lambda \times N_d},
\end{equation}
indexed by delay $\tau$, probe coordinate $\lambda$ (wavelength, wavenumber or
energy) and detector $d$. The accompanying axes are the delay vector
$\tau \in \mathbb{R}^{N_\tau}$ and the probe vector
$\lambda \in \mathbb{R}^{N_\lambda}$. Table~\ref{tab:fields} lists the seven
fields every loader returns.

\begin{table}[h]
  \centering
  \begin{tabular}{lll}
    \toprule
    Field & Symbol / shape & Meaning \\
    \midrule
    \code{Zavg\_R} & $N_\tau \times N_\lambda \times N_d$ & averaged signal \\
    \code{delays}  & $N_\tau$ & delay axis \\
    \code{probe}   & $N_\lambda$ & probe (spectral) axis \\
    \code{Units}   & dict & unit labels (see \S\ref{sec:settings}) \\
    \code{Nscans}  & scalar & number of scans (\code{nan} if unknown) \\
    \code{Zss\_R}  & $N_\tau \times N_\lambda \times N_d \times N_s$ & single-scan signal \\
    \code{Zstdv}   & $N_\tau \times N_\lambda$ & standard deviation \\
    \bottomrule
  \end{tabular}
  \caption{Fields returned by a 1-D loader and stored on \code{Dataset1D}.}
  \label{tab:fields}
\end{table}

The most-processed signal is exposed through the \code{Z} property, which
returns the background-corrected array once a correction has been applied and
falls back to the raw array otherwise.

\subsection{Loading and the loader registry}
Data formats are decoupled from the dataset through a \emph{loader registry}
(\code{pymorgan.oneD.load}). A loader is any callable that maps a file path to
the seven fields of Table~\ref{tab:fields}. Loaders are registered under one or
more data-type names and resolved at load time:

\begin{lstlisting}
import pymorgan as pm

data = pm.load_1D("scan.pdat", data_type="PDAT")
print(pm.available_loaders())
\end{lstlisting}

The currently registered names are \code{PDAT} (implemented) together with the
instrument-specific formats \code{UniGE\_fsTA}, \code{UniGE\_nsTA},
\code{MESS\_TRIR}, \code{MESS\_TRUVIS}, \code{UniGE\_FLUPSold},
\code{UniGE\_FLUPSnew}, \code{Helios\_TA} and \code{UoS\_IRpp}, which are
recognised but raise \code{NotImplementedError} until their on-disk formats are
defined. A new format is added by decorating a reader:

\begin{lstlisting}
@pm.register_loader("MyInstrument")
def read_my_instrument(path):
    ...
    return Zavg_R, delays, probe, Units, Nscans, Zss_R, Zstdv
\end{lstlisting}

\paragraph{The PDAT format.}
A \code{.pdat} file is comma-separated. The first line carries the time and
probe units separated by an asterisk (e.g.\ \code{ps*cm\textasciicircum\{-1\}}).
The first data row holds the probe axis (with a leading $0$), and each
subsequent row begins with a delay value followed by the signal at every probe
coordinate. PDAT files store averaged data already in m\,OD, so no rescaling is
applied on read.

\subsection{Background correction}
The pre-time-zero background is estimated by averaging the signal over a delay
window $[\tau_\mathrm{min},\tau_\mathrm{max}]$ and subtracting it from every
delay,
\begin{equation}
  Z_C(\tau,\lambda,d) = Z(\tau,\lambda,d)
    - \big\langle Z(\tau',\lambda,d) \big\rangle_{\tau' \in [\tau_\mathrm{min},\tau_\mathrm{max}]},
\end{equation}
where the average is taken with \code{nan}-aware statistics
(\code{pymorgan.oneD.process}). The operation is applied through

\begin{lstlisting}
data.background_correct(tmin=-20, tmax=-5)          # apply
data.background_correct(tmin=-20, tmax=-5, do_correct=False)  # passthrough
\end{lstlisting}

With \code{do\_correct=False} the raw signal is copied unchanged into the
corrected fields, so downstream code can always read \code{Z} uniformly. The
method returns \code{self}, allowing calls to be chained.

\subsection{Solvent subtraction}\label{sec:solvent}
For measurements in solution, the pure-solvent response often has to be
subtracted. The module exposes two entry points.

\paragraph{Manual subtraction.} A scaled, time-shifted solvent trace is removed
from the sample signal:
\begin{lstlisting}
data.subtract_solvent(solvent, dt=0.0, scale=1.0)
\end{lstlisting}
\code{dt} shifts the solvent delay axis before interpolation (accounts for a
timing offset between the two measurements) and \code{scale} multiplies the
solvent signal before subtraction.

\paragraph{Automatic solvent fitting.} \code{fit\_solvent\_auto} models the
solvent contribution pixel by pixel using the solvent trace convolved with the
Instrument Response Function (IRF) and removes it:
\begin{lstlisting}
fit = data.fit_solvent_auto(solvent)
data.subtract_solvent_fit(fit)
\end{lstlisting}
Optional per-pixel fits for the time-zero offset and the amplitude scale can be
enabled through the \code{per\_pixel\_dt} and \code{per\_pixel\_scale} settings
(Section~\ref{sec:settings}); the defaults are \code{false} and \code{true},
respectively.

\subsection{Shockwave subtraction}\label{sec:shockwave}
In transient IR / UV-Vis measurements using flowing samples, acoustic shockwaves
generated by the pump beam can appear as a time-dependent artefact that is
correlated across \emph{all} probe wavelengths. The artefact is removed by
averaging a selected subset of detector pixels (or a probe-axis range) into a
one-dimensional shockwave trace and subtracting it from every pixel:
\begin{lstlisting}
from pymorgan.oneD.process import subtract_shockwave

Z_corrected = subtract_shockwave(Z, pixels=[0, 1, 2])   # pixel indices (0-based)
\end{lstlisting}

The GUI exposes this through the \textbf{Subtract Shockwave\ldots} button in the
1-D tab, which opens the \code{ShockwaveSubtractionDialog}. The dialog lets the
user choose between pixel-index and probe-axis-value selection, displays a live
preview of the proposed shockwave trace, and applies the correction on
confirmation (see \S\ref{sec:oneD:gui:dialogs}).

\subsection{Chirp (dispersion) correction}\label{sec:chirp}
Ultrafast time-resolved spectroscopy measurements often suffer from Group Delay Dispersion (GDD), commonly referred to as ``chirp''. As the probe pulse passes through various dispersive optical elements (such as white-light generation crystals, lenses, sample cell windows, and the solvent itself), different spectral components are temporally stretched. Consequently, the effective time-zero $t_0(\lambda)$ varies across the probe wavelength range $\lambda$.

PyMORGAN provides tools to characterise and correct this dispersion. The time-zero profile $t_0(\lambda)$ is modelled using a generalised Cauchy dispersion equation:
\begin{equation}
  t_0(\lambda) = c_0 + \sum_{n=1}^{N_C-1} c_n \left(\frac{\lambda_{\mathrm{ref}}}{\lambda}\right)^{2n},
\end{equation}
where $c_0, \dots, c_{N_C-1}$ are the dispersion coefficients, $\lambda_{\mathrm{ref}}$ is a reference wavelength (which defaults to the mean of the fitted range to make the coefficients dimensionless and of comparable magnitude), and $N_C$ is the number of terms (usually $6$).

To determine $t_0(\lambda)$ across the probe range, PyMORGAN implements three fit modes in the \code{pymorgan.oneD.chirp} module:
\begin{description}
  \item[Automatic] Fits the coherent artefact and early dynamics at every pixel independently. The model consists of a Gaussian Instrument Response Function (IRF), its optional first and second derivatives, a constant offset, and an optional erf-broadened exponential decay. PyMORGAN uses Variable Projection (VARPRO) to eliminate the linear amplitudes in closed form, optimizing only the nonlinear parameters $(t_0, \mathrm{FWHM}, \tau_{\mathrm{exp}})$ per pixel. A second pass uses smoothed parameters from the first pass as initial guesses to improve stability.
  \item[Step function] A simplified alternative that jointly fits a Gaussian IRF (+ derivatives) and a Heaviside step convolved with the same Gaussian over a narrow early-time window. This is robust when the early-time exponential decay is not needed or is ill-conditioned.
  \item[Manual] Fits the Cauchy dispersion equation directly to a set of user-picked $(\lambda, t_0)$ points (e.g., clicked interactively on the contour map in the GUI), bypassing the per-pixel IRF fitting.
\end{description}

Once $t_0(\lambda)$ is obtained, the Cauchy dispersion equation is fit using bound-constrained, iteratively reweighted linear least-squares (IRLS with Tukey's bisquare weights) to handle outliers and reject bad pixel fits.

To apply the correction, the dataset's delay axes are interpolated onto a common chirp-free delay grid. Two boundary handling strategies are available:
\begin{itemize}
  \item \code{crop}: Slices the delay axis to keep only the delay region where all wavelength channels are defined, discarding the out-of-bounds regions that become \code{NaN} after shifting.
  \item \code{extrapolate}: Fills out-of-bound values with the nearest valid boundary endpoint value, keeping the original delay range.
\end{itemize}

\begin{lstlisting}
# Fit chirp automatically and apply the correction
fit = data.fit_chirp_automatic(deriv_mask=(True, True), use_exp=True)
data.plot_chirp_diagnostics(fit)
data.apply_chirp_correction(fit, boundary_handling="crop")
\end{lstlisting}

\subsection{Plotting}
Five plotters live in \code{pymorgan.oneD.plot}; each is exposed both as a
\code{Dataset1D} method and as a stage-level function that takes the dataset as
its first argument, so a new view only needs \code{(data, ax, \dots)} and
inherits the whole pipeline's context. Each returns the matplotlib axis (the
kinetic view also returns the extracted traces). They are listed in
Table~\ref{tab:plotters}.

\begin{table}[h]
  \centering
  \begin{tabular}{ll}
    \toprule
    Method / function & View \\
    \midrule
    \code{plot\_contour}        & delay-vs-probe contour map \\
    \code{plot\_surface}        & filled delay-vs-probe surface \\
    \code{plot\_spectra}        & transient spectra at selected delays \\
    \code{plot\_kinetics}       & kinetic traces at selected probe positions \\
    \code{plot\_species\_spectra} & EAS/SAS from a global-analysis fit \\
    \bottomrule
  \end{tabular}
  \caption{The 1-D plotters (\code{pymorgan.oneD.plot}).}
  \label{tab:plotters}
\end{table}

\begin{lstlisting}
data.plot_contour(Zscale=20)                          # method form
data.plot_spectra([0.5, 1, 5, 20, 100], doSmooth=1)
ax, t, Y = data.plot_kinetics([2132, 2218], plotStyle="-")

from pymorgan.oneD import plot                          # stage-function form
plot.plot_contour(data, label_style="/", cmap_ID="Rd/Wh/Bu v2")
\end{lstlisting}

The label style, colourmap, time-axis scale and \dA{} unit convention are taken
from the active settings (\S\ref{sec:settings}); any of them can be overridden
per figure through the corresponding keyword argument.

For datasets with high dynamic range where weak features would otherwise be washed out by dominant bleach or ESA peaks, \code{plot_contour} provides non-linear $\operatorname{arcsinh}$ colour scaling via \code{Asinh=True} (also toggled in the GUI plot controls). Rather than modifying the underlying signal values, PyMORGAN applies an inverse hyperbolic sine colour normalization (\code{matplotlib.colors.AsinhNorm}) paired with non-linear $\sinh$-spaced contour levels:
\begin{equation}
  z_k = s \cdot \sinh(u_k), \quad u_k \in \left[-\operatorname{arcsinh}\left(\frac{|Z_{\min}|}{s}\right), \operatorname{arcsinh}\left(\frac{Z_{\max}}{s}\right)\right],
\end{equation}
where $s = \frac{p}{100} \cdot \max(|Z_{\min}|, |Z_{\max}|)$ is the linear crossover threshold set by \code{asinh_pct} ($p$, default $5.0\%$). Signals within $\pm s$ behave linearly, while larger amplitudes are compressed logarithmically. Because the underlying data array is not mutated, coordinates and hover values retain their physical units ($\mathrm{mOD}$), and colorbar ticks remain linear and clean (centred at zero and rounded to sensible numbers).

\subsection{Composing figures and displaying}
By default every plotter creates its own figure and axis. Passing an existing
axis through \code{ax=} draws into it instead, which is how panels are combined
into a layout or several datasets are overlaid on one axis. Repeated elements
(axis labels and, for the maps, the colorbar) can be switched off so they are
not drawn more than once; the toggles are listed in Table~\ref{tab:toggles}.

\begin{table}[h]
  \centering
  \begin{tabular}{lll}
    \toprule
    Keyword & Default & Plotters \\
    \midrule
    \code{show\_xlabel}        & \code{True} & all \\
    \code{show\_ylabel}        & \code{True} & all \\
    \code{show\_colorbar}      & \code{True} & \code{plot\_contour}, \code{plot\_surface} \\
    \code{show\_colorbar\_label} & \code{True} & \code{plot\_contour}, \code{plot\_surface} \\
    \bottomrule
  \end{tabular}
  \caption{Per-call toggles for composing multi-panel figures.}
  \label{tab:toggles}
\end{table}

\begin{lstlisting}
import matplotlib.pyplot as plt
fig, (axL, axR) = plt.subplots(1, 2)

# two maps side by side; suppress the duplicated colorbar / y-axis label
dataA.plot_contour(ax=axL, show_colorbar=False)
dataB.plot_contour(ax=axR, show_ylabel=False)

# overlay two datasets' transient spectra on a single axis
dataA.plot_spectra([1, 5, 20], ax=axL)
dataB.plot_spectra([1, 5, 20], ax=axL, show_ylabel=False)

pm.show_plots()   # wrapper around plt.show(); pm.show is an alias
\end{lstlisting}

When a plotter is given an existing axis it does not resize the parent figure,
so dropping a plot into a layout leaves the surrounding panels untouched.
\code{pm.show\_plots(block=False)} returns immediately (e.g.\ inside a GUI event
loop), and \code{pm.close\_plots()} closes all open figures.

\subsection{Kinetic analysis}\label{sec:analyse}
Kinetic analysis is \emph{not} part of PyMORGAN. Trace extraction and fitting
--- multi-exponential decays, and full global/target analysis with rate-matrix
models (the source of the EAS/SAS consumed by \code{plot\_species\_spectra})
--- live in the companion project \textbf{PyRATE-TA}, which shares this
loader/dataset layer.

PyMORGAN keeps the two ends of that seam. Traces are obtained from
\code{plot\_kinetics}, whose third return value is the extracted data itself:

\begin{lstlisting}
ax, t, Y = data.plot_kinetics([1950, 2000], plotStyle="-")
Y.shape                                           # [Ndelays x Ncuts]
\end{lstlisting}

and fitted spectra are rendered by \code{plot\_species\_spectra}, which takes
plain arrays (\code{Sfit}, \code{Taus}, \code{TauErr}, \code{isFixTau},
\code{modelType}) and is therefore independent of how the fit was produced.

The third piece of the seam is the noise. \code{Dataset1D.noise\_array()}
returns the per-point standard deviation --- from \code{Zstdv} when a sibling
\code{.pdatn} file was loaded, otherwise the spread across single scans, and
\code{None} when the dataset carries neither:

\begin{lstlisting}
sigma = data.noise_array()      # [Ndelays x Npixels x Ndetectors], or None
\end{lstlisting}

A weighted fit uses it as \code{1/sigma}, which is what turns a fit statistic
into a true reduced $\chi^2$; without it the statistic is a plain sum of
squared residuals. PyMORGAN owns the noise because it owns the loaders; PyRATE-TA
never estimates it.

\subsection{Settings and aesthetics}\label{sec:settings}
Presentation choices that are not intrinsic to the data are collected in a
single \code{Settings} object. It is framework-agnostic: the GUI binds to the
fields it describes, but the package never depends on the GUI. The fields are
summarised in Table~\ref{tab:settings}.

\begin{table}[h]
  \centering
  \begin{tabular}{lll}
    \toprule
    Field & Values & Effect \\
    \midrule
    \code{profile} & regular / poster / inset & selects a \code{.mplstyle} file \\
    \code{font\_scale} & float & multiplies the profile font sizes \\
    \code{label\_style} & \code{()} \code{[]} \code{/} (empty) & axis-label delimiter \\
    \code{delta\_a\_units} & \code{mOD} / \code{x1E3} & \dA{} label convention \\
    \code{cmap} & string & default contour colourmap \\
    \code{time\_axis\_scale} & symlog / log / lin & default time-axis scale \\
    \code{freq\_label} & Omega\_n / Omega\_PP / \dots & 2-D frequency-axis convention \\
    \code{x\_axis\_unit} & native / nm / cm-1 / eV / THz & spectral X-axis display unit \\
    \code{secondary\_axis} & bool & complementary X axis (top) \\
    \code{chirp\_boundary\_handling} & crop / extrapolate & boundary handling for chirp correction \\
    \bottomrule
  \end{tabular}
  \caption{Fields of the \code{Settings} object.}
  \label{tab:settings}
\end{table}

The \code{label\_style} field controls the delimiter so that an axis reads, for
example, ``Wavelength (nm)'', ``Wavelength [nm]'' or ``Wavelength / nm''. The
\code{delta\_a\_units} field switches the signal-axis label between
``\dA{} / m\,OD'' and ``\dA{} $\times 10^{3}$''; because $1\,\mathrm{mOD}
\equiv \dA \times 10^{3}$, this is a labelling choice only and never rescales
the stored data. The active \code{label\_style} delimiter is honoured for the
m\,OD label on every axis, including the contour colorbar (so ``/'' yields
``\dA{} / m\,OD''); the $\times 10^{3}$ form is always shown without a delimiter.

The style profiles wrap the existing hand-tuned \code{.mplstyle} files rather
than replacing them; \code{font\_scale} is applied as a light multiplicative
overlay on top. Text is always rendered with matplotlib's mathtext engine; the
LaTeX backend (\code{text.usetex}) is never enabled.

The \code{x\_axis\_unit} field converts the spectral (probe) X axis of the transient-spectra and contour plots between wavelength (nm), wavenumber (cm$^{-1}$), energy (eV) and frequency (THz); \code{native} keeps each dataset's stored unit. The conversion is physical (the axis values are recomputed, pivoting through the wavelength), not a relabelling. When \code{secondary\_axis} is enabled a second X axis is drawn on top in the \emph{complementary} unit: cm$^{-1}\leftrightarrow$nm, with nm used as the complement for eV and THz. Both fields default to the active settings and can be overridden per call, e.g. \code{data.plot\_spectra([1, 5, 20], x\_axis\_unit="nm", secondary\_axis=True)}. The earlier \code{dualScale} keyword is retained as a deprecated alias (a truthy value maps to \code{secondary\_axis=True}). Whenever a wavenumber axis (main or secondary) exceeds $5000\,\mathrm{cm^{-1}}$ its tick labels are divided by $1000$ and the axis is relabelled ``Wavenumber ($\times 10^{3}$ cm$^{-1}$)'', so UV/Vis wavenumbers stay legible; IR ranges are left unscaled. The axis data, stored limits and GUI spin boxes remain in true cm$^{-1}$ -- only the displayed tick labels and unit are scaled.

The \code{chirp\_boundary\_handling} field controls how out-of-bounds delays are treated during the chirp correction interpolation step, selecting between \code{crop} and \code{extrapolate} (see \S\ref{sec:chirp}).

Settings are read from and written to a TOML file, with comments preserved on
write so the file remains hand-editable after the GUI saves it:

\begin{lstlisting}
pm.load_settings("settings.toml")   # make the file's settings active
pm.update_settings(profile="poster", label_style="/")
pm.apply_style()                    # apply the profile to matplotlib
pm.save_settings("settings.toml")   # persist (keeps comments)
\end{lstlisting}

A module-level \emph{active} settings object is read by the plotting methods.
Temporary overrides can be scoped with a context manager:

\begin{lstlisting}
with pm.use_settings(delta_a_units="mOD"):
    data.plot_spectra([1, 5, 20])   # this figure only
\end{lstlisting}

\subsection{Worked example}
A complete script reproducing the standard workflow is provided in
\code{examples/pump\_probe\_1D.py}. In outline, the four stages read:

\begin{lstlisting}
import pymorgan as pm

pm.load_settings("settings.toml")
pm.apply_style()

data = pm.load_1D("testData/testTRIR.pdat", data_type="PDAT")  # load
data.background_correct(tmin=-20, tmax=-5)                       # process
data.plot_contour(Zscale=20, Yscale="lin")                      # plot
ax, t, Y = data.plot_kinetics([1980, 2030], plotStyle="-")      # plot + extract
\end{lstlisting}

\subsection{Graphical interface — 1-D dialogs}\label{sec:oneD:gui:dialogs}

Several interactive dialogs supplement the main 1-D tab.

\subsubsection{Chirp correction dialog}

The \textbf{Chirp\ldots} button opens \code{chirp\_dialog.py}, which guides
the user through chirp correction in three sub-panels:
\begin{enumerate}
  \item \textbf{Mode selection} --- choose Automatic, Step-function or Manual.
  \item \textbf{Parameter input} --- set the delay window, IRF FWHM, and
        (for Automatic mode) which derivative terms to include.
  \item \textbf{Diagnostics} --- the fitted $t_0(\lambda)$ curve overlaid on
        the contour map; the boundary-handling strategy (crop / extrapolate)
        is chosen here before the correction is applied.
\end{enumerate}
A progress dialog (\code{chirp\_progress\_dialog.py}) reports per-pixel fit
progress for the Automatic mode, which can be parallelised via the
\code{parallel\_fitting} setting (thread pool, process pool, or sequential).

\subsubsection{Solvent subtraction dialog}

The \textbf{Subtract Solvent\ldots} button opens a dialog in
\code{dialogs.py} that loads a pure-solvent scan, previews the subtraction,
and applies \code{fit\_solvent\_auto} or the manual \code{subtract\_solvent}
path according to the chosen settings.

\subsubsection{Shockwave subtraction dialog}

The \textbf{Subtract Shockwave\ldots} button opens
\code{shockwave\_dialog.py}. The user selects a set of pixels (by index or
probe value) whose average forms the shockwave reference trace, inspects a
live preview, and confirms. The correction calls \code{subtract\_shockwave}
and updates the dataset in place.
